% TOP_PROBLEM: {"id": 30003897, "title": "Green’s Canonical Syzygy Conjecture for Arbitrary Curves", "category_id": 5, "category": {"id": 5, "name": "algebraic_geometry", "display_name": "Algebraic Geometry", "description": "Geometric objects defined by polynomial equations.", "slug": "algebraic-geometry", "order_index": 5, "created_at": "2026-07-31T15:26:25.671Z"}, "status": "open"}
% FIRST_POSTED: 2026-09-25
% ENTRY_KIND: statement_only
% !TeX program = lualatex
% Definition and sources only; this file is not an LLM solution attempt.
\documentclass[11pt]{article}
\usepackage[margin=25mm]{geometry}
\usepackage{fontspec}
\setmainfont{Latin Modern Roman}
\usepackage{amsmath,amssymb,mathtools}
\usepackage{unicode-math}
\setmathfont{Latin Modern Math}
\usepackage{xurl,hyperref}
\hypersetup{colorlinks=true,urlcolor=blue}
\setcounter{secnumdepth}{0}
\setlength{\parindent}{0pt}
\setlength{\parskip}{5pt}
\setlength{\emergencystretch}{3em}
\begin{document}
\section*{\texorpdfstring{Green’s Canonical Syzygy Conjecture for Arbitrary Curves}{Green’s Canonical Syzygy Conjecture for Arbitrary Curves}}\label{30003897}
\subsection{Definitions and mathematical statement}
In the complex characteristic-zero setting of the cited source, let $C$ be a smooth nonhyperelliptic curve of genus $g$, $V=H^0(C,K_C)$, and define $K_{p,1}(C,K_C)$ as the middle cohomology of the Koszul complex
\[
 \bigwedge^{p+1}V\otimes H^0(C,\mathcal O_C)
 \longrightarrow\bigwedge^pV\otimes H^0(C,K_C)
 \longrightarrow\bigwedge^{p-1}V\otimes H^0(C,K_C^2).
\]
The maps alternate multiplication of sections. The Clifford index is the minimum of $\deg L-2(h^0(L)-1)$ over line bundles with $h^0(L),h^1(L)\ge2$, with the standard value one for a nonhyperelliptic genus-three curve. The selected Green equivalence is
\[
 [K_{\ell,1}(C,K_C)=0\text{ for every }\ell\ge p]
 \quad\Longleftrightarrow\quad\operatorname{Cliff}(C)>g-p-2.
\]
The catalogue does not name a ground field; the source setting is stated explicitly rather than extending this assertion silently to arbitrary positive characteristic.
\par\smallskip\textit{Formulation and concept references:} \href{https://webusers.imj-prg.fr/~claire.voisin/Articlesweb/syzod.pdf}{[S1]}.

\subsection{Short English statement}
For every smooth nonhyperelliptic complex curve, does its Clifford index determine exactly where the linear syzygies of its canonical embedding vanish?
\subsection{Sources}
\begin{itemize}
\item[S1]Claire Voisin, Green's canonical syzygy conjecture for generic curves of odd genus, Compositio Mathematica, 2005.
\url{https://webusers.imj-prg.fr/~claire.voisin/Articlesweb/syzod.pdf}.
\textit{Formulation source.}
\end{itemize}
\end{document}
